% =====================================================================
%  POST-ACCEPTANCE typeset version -- De Gruyter / JQAS house style.
%
%  This is NOT the file to submit. JQAS reviews blind, on 8.5x11 double-
%  spaced pages; that copy is main.tex. This one uses the publisher's
%  dgruyter.sty to preview how the article sets in the journal, and it
%  carries the author identity that main.tex deliberately omits.
%
%  Both documents \input the same body.tex, so the prose cannot drift
%  between them -- only the front matter and page geometry differ.
%
%  Build:  make paper-dgruyter
%
%  Before sending to production, fill in the \received/\revised/\accepted
%  dates, volume, issue, and DOI supplied by the editorial office, and
%  restore the real author block below.
% =====================================================================
\documentclass[USenglish]{article}

\usepackage[utf8]{inputenc}
\usepackage[big,online]{dgruyter}
\usepackage{lmodern}
\usepackage{microtype}
\usepackage[round,authoryear]{natbib}

\usepackage{booktabs}
\usepackage{graphicx}
\usepackage{amsmath}
\usepackage{amssymb}

\bibliographystyle{agsm}   % Harvard author-date, as JQAS requires

% Macros shared with main.tex via body.tex.
% Macros used by body.tex, which both main.tex (submission) and
% main-dgruyter.tex (post-acceptance typeset) \input. Anything the shared body
% depends on belongs here, NOT in a single document's preamble -- otherwise the
% other document fails with "undefined control sequence".
%
% \crvfigure is deliberately NOT defined here: each document needs its own
% version (placement marker vs. embedded graphic), so they define it themselves.

% Long snake_case identifiers set in \texttt are single unbreakable words and
% punch through the margin. \crvus prints an underscore and permits a line
% break after it, so such names wrap at a natural seam. Call it as \crvus{} so
% the following space is not swallowed by the macro name.
\newcommand{\crvus}{\_\allowbreak}

% Double-spaced 12pt on a single-column measure gives TeX few good break points,
% and paragraphs dense with hyphenated compounds ("five-module pipeline---
% ingestion, ... alternate-state generation") exceed \tolerance and set overfull
% rather than loose. \emergencystretch authorizes a final pass with extra
% interword stretch, which resolves those without touching the prose.
\setlength{\emergencystretch}{2em}


% The typeset version embeds the artwork -- the "separate files" rule is a
% submission requirement, not a typesetting one. Same macro signature as
% main.tex so body.tex is byte-identical between the two documents.
\newcommand{\crvfigure}[3]{%
  \begin{figure}[htbp]
    \centering
    \includegraphics[width=#3\textwidth]{../outputs/figures/#2.pdf}
    \caption{See Figure Legends.}
    \label{#1}
  \end{figure}}

\begin{document}

%%%--------------------------------------------%%%
  \articletype{Research Article}
  \received{Month DD, YYYY}
  \revised{Month DD, YYYY}
  \accepted{Month DD, YYYY}
  \journalname{Journal of Quantitative Analysis in Sports}
  \journalyear{YYYY}
  \journalvolume{XX}
  \journalissue{X}
  \startpage{1}
  \aop
  \DOI{10.1515/jqas-YYYY-XXXX}
%%%--------------------------------------------%%%

\title{Challenge Run Value ($cRV$): Quantifying Strategic Value in MLB's ABS
Challenge Era}
\runningtitle{Challenge Run Value in MLB's ABS era}

% TODO(post-acceptance): restore the real author block. Left as placeholders
% so this file can never leak identity if it is compiled by mistake.
\author*[1]{Author Name}
\runningauthor{A.~Name}
\affil[1]{\protect\raggedright
Institution, Department, City, Country, e-mail: author@example.edu}

\abstract{We introduce Challenge Run Value ($cRV$), a pitch-level metric for
valuing Automated Ball-Strike (ABS) challenges in Major League Baseball. The
model combines run-expectancy deltas from count- and plate-appearance-ending
state changes with a time-decaying opportunity-cost penalty for failed
challenges. Unlike framing metrics, which measure physical receiving skill
under human umpiring, $cRV$ treats the challenge as a scarce strategic
resource whose value depends on timing, conversion, and forgone future
options.

We implement a reproducible five-module pipeline---ingestion, event isolation,
alternate-state generation, RE288 mapping, and scoring---with bootstrap
uncertainty, leverage stratification, and sensitivity analysis. A central
methodological contribution is the valuation of terminal outcomes: challenges
converting a called strike three into ball four, or the reverse, end the plate
appearance and must be valued using post-event base/out states and immediate
runs scored rather than zeroed placeholders. Such events are the
highest-leverage challenges in the game, and zeroing them erases the metric's
most important signal, inverting the sign of the aggregate result.

Validation on a designed corpus of ten challenge events yields mean
$cRV = 0.193$ runs per event, with the largest contributions from a
bases-loaded walk conversion and a late-inning inning-ending strikeout
conversion. Full-season estimates remain the primary empirical target; current
figures are a methodological baseline, not a player ranking.}

\keywords{run expectancy; opportunity cost; catcher framing; sabermetrics;
decision quality; resource allocation}

\maketitle

\section{Introduction}
MLB's adoption of the Automated Ball-Strike (ABS) challenge system reframes
pitch evaluation. Under the framing-era paradigm, catcher value at the margin
of the strike zone was largely a function of receiving technique and umpire
interaction \citep{marchi2013,fast2011,mlbframing}. A catcher earned runs by
making borderline pitches \emph{look} like strikes, and the skill was
continuous, repeatable, and accumulated over thousands of receptions per
season.

The ABS challenge system replaces that continuous skill with a discrete,
scarce resource. Each team carries a limited number of challenges; a
successful challenge is retained, a failed challenge is consumed. The
strategic problem is therefore no longer ``how do I make this pitch look like
a strike'' but ``is this pitch worth spending a challenge on, given how much
game remains and how likely I am to be right.'' That is a resource-allocation
problem over a nine-inning (or longer) horizon, and it is not measured by any
existing public metric.

Existing public ABS coverage has been largely descriptive: challenge rates,
success percentages, and umpire-agreement rates \citep{mlbabs2026}. These are
useful accuracy diagnostics but they are denominated in \emph{challenges}, not
runs, and they are indifferent to context. A catcher who wins 60\% of his
challenges on 0--0 counts in the first inning and a catcher who wins 60\% of
his challenges on full counts with the bases loaded in the ninth have
identical success rates and vastly different value.

What has been missing is a unified, run-denominated valuation that (i)~credits
the immediate state change produced by an overturn and (ii)~charges the
opportunity cost of a failed challenge against the remaining game horizon.
This paper formalizes that valuation as Challenge Run Value ($cRV$), implements
it as an open and reproducible pipeline, validates the implementation against a
designed corpus of edge cases, and applies it to the 2026 regular season
through August 12: $1{,}971$ plate-appearance-ending challenge events scored
against an empirically estimated RE288 surface.

\subsection{Contributions}
\begin{enumerate}
  \item A formal run-denominated metric, $cRV$, combining an RE288 pitch-state
        delta with a time-decaying opportunity-cost penalty
        (Section~\ref{sec:method}).
  \item Correct treatment of plate-appearance-ending challenges: walk force
        chains with runs scored, and strikeouts with out increments and
        inning-end zeroing (Section~\ref{sec:terminal}). We show this is not a
        detail: it is where nearly all of the metric's signal lives.
  \item A reproducible five-module pipeline with provenance tracking, bootstrap
        uncertainty, leverage stratification, an empirical $EV_c$ estimate, and
        a win-denominated overlay (Section~\ref{sec:pipeline}), backed by a
        designed validation corpus and unit-test suite that exercises each
        state-transition path (Section~\ref{sec:validation}).
  \item An empirical application to the 2026 regular season through
        August 12 ($1{,}971$ plate-appearance-ending challenges scored against
        an empirically estimated RE288 matrix) that
        quantifies the run value of the challenge as used in practice and shows
        catchers extract roughly two-and-a-half times the per-challenge value
        of batters (Section~\ref{sec:results}).
\end{enumerate}

\subsection{Related work}
Count-dependent run expectancy is foundational in sabermetric modeling
\citep{tango2007}. The classical RE24 matrix conditions expected runs on the 24
base-out states. The RE288 extension (8 base states $\times$ 3 outs $\times$
12 counts) conditions additionally on the ball-strike count, which is
precisely the dimension an ABS challenge manipulates
\citep{tangotiger2018}. Without the count dimension, a challenge that
moves a count from 0--1 to 1--0 is invisible; RE288 makes it measurable.

Catcher framing research established that called-strike environments are
measurable and strategically relevant \citep{marchi2013,fast2011,mlbframing}.
That literature is the natural predecessor to this work, and also its natural
contrast: framing measured \emph{how often} borderline pitches were converted
to strikes through receiving skill, whereas $cRV$ measures \emph{when} a
challenge should be spent and \emph{whether} spending it paid off. Framing runs
and $cRV$ are complementary rather than competing: in a full-ABS world the
former decays toward zero while the latter becomes the live skill.

The strategic problem of \emph{when} to spend a scarce, retained-on-success
review resource has an established economics and operations-research
literature. \citet{abramitzky2012} characterize optimal
Hawk-Eye line-call challenge behavior in professional tennis; \citet{clarkenorman2012} solve the challenge decision as a dynamic
program; and \citet{shivakumar2018} documents strategic review use
in cricket's Decision Review System, whose ``retained on success'' rule
mirrors ABS. Broader work on umpiring-context effects and decision incentives
in sport \citep{scorecasting} frames the same problem at the behavioral level.
Our opportunity-cost penalty is a deliberately simple approximation to the
option value these dynamic-program treatments compute exactly; Section
\ref{sec:limitations} returns to the gap between the two.

\section{Materials and Methods}

\subsection{Methodology}
\label{sec:method}

\subsubsection{Notation}
For challenge event $i$, let the \emph{pre-pitch state} be the triple
$(c_i, b_i, o_i)$ where $c_i$ is the ball-strike count, $b_i \in \{0,1\}^3$ is
the base occupancy string (1B, 2B, 3B), and $o_i \in \{0,1,2\}$ is the out
count. Let $RE(\cdot)$ denote the run value of a post-pitch state, defined to
include both any runs scored immediately on the play and the run expectancy of
the resulting state.

Every challenge concerns a called pitch that the umpire (or ABS system) ruled
either a ball or a strike. Two realities follow:
\begin{itemize}
  \item the \textbf{original} reality, in which the original call stands, and
  \item the \textbf{overturned} reality, in which the call is reversed.
\end{itemize}
If the challenge fails, the overturned reality never materializes and the two
coincide.

\subsubsection{Pitch-state delta}
\begin{align}
\Delta RV_{pitch,i} &=
\begin{cases}
RE(\text{overturned}_i) - RE(\text{original}_i)
  & \text{batter challenge succeeds}, \\[4pt]
RE(\text{original}_i) - RE(\text{overturned}_i)
  & \text{catcher challenge succeeds}, \\[4pt]
0 & \text{challenge fails}.
\end{cases}
\end{align}

The sign convention encodes whose interest the challenge serves. Run
expectancy is always stated from the batting team's perspective, so a batter
benefits when the overturn \emph{raises} RE (a strike becomes a ball) and a
catcher benefits when the overturn \emph{lowers} it (a ball becomes a strike).
Taking the catcher's delta as $RE_{\text{original}} - RE_{\text{overturned}}$
makes a successful defensive challenge score positive, so that $cRV$ is
directly comparable across roles.

Failed challenges leave the on-field state unchanged. The pitch delta is
exactly zero (not merely small) because the umpire's call stands and play
proceeds from the original reality.

\subsubsection{Opportunity-cost penalty}
Let $EV_c$ be the baseline expected value of \emph{retaining} a challenge. In
v1 we fix $EV_c = 0.15$ runs as a transparent constant; Section
\ref{sec:sensitivity} reports sensitivity over $[0.034, 0.30]$. Conceptually,
\begin{align}
EV_c = P(\text{future need}) \times P(\text{win}) \times
       \overline{\Delta RV}_{\text{late}},
\end{align}
the probability a challenge will be wanted again, times the conversion rate,
times the average value of a late-game conversion. Estimating the factors from
the 2026 corpus gives $P(\text{future need}) = 0.206$ (the share of challenges
that are not the challenging team's last challenge of the game),
$P(\text{win}) = 0.465$, and $\overline{\Delta RV}_{\text{late}} = 0.353$
runs, for an empirical $EV_c = 0.034$ runs (cluster-bootstrap 95\% CI
$[0.030,\,0.038]$; Table~\ref{tab:evc}). This is a calibration of the penalty
constant rather than a direct measurement of option value: $P(\text{future
need})$ is a within-game frequency proxy, not a model of future game states.
Because the empirical value sits below the $0.15$ baseline, the headline
results that follow are computed under a deliberately conservative $EV_c$.

\begin{table}[h]
\centering
\caption{Empirical estimate of the retained-challenge value $EV_c$.}
\label{tab:evc}
\small
\begin{tabular}{lrrr}
\toprule
Metric & Value & CI lower & CI upper \\
\midrule
P(future need) & $0.206$ & $0.191$ & $0.222$ \\
P(win) & $0.465$ & $0.444$ & $0.486$ \\
Mean late-game delta RV & $0.353$ & $0.326$ & $0.384$ \\
Empirical EV\_c & $0.034$ & $0.030$ & $0.038$ \\
\bottomrule
\end{tabular}
\end{table}


Failed challenges incur
\begin{align}
Penalty_i = -\left(\frac{\max(0,\, 9 - \mathrm{inning}_i)}{9}\right)
            \frac{EV_c}{C_i},
\label{eq:penalty}
\end{align}
where $C_i$ is the number of challenges the team held immediately before the
failure. The penalty prices the opportunity cost of the challenge the failure
destroys: each team enters the game with a two-challenge budget
\citep{mlbabs2026} (a failure consumes one, a success retains it), so losing
the \emph{last} challenge forfeits the entire retained option value, while
losing one of two forfeits only a $1/C$ share. In the 2026 corpus, $973$ of
the $1{,}055$ failed challenges were the team's first failure of the game
($C_i = 2$, a half charge) and $82$ were its last ($C_i = 1$, a full charge).

Four properties of Equation~\ref{eq:penalty} deserve comment. First,
successful challenges are retained under ABS rules \citep{mlbabs2026}, so $Penalty_i = 0$ on
wins, since the resource was not consumed. Second, the penalty conditions on the
actual scarce resource: a failure with a reserve challenge still in hand
($C_i = 2$) is charged half of one with none in reserve ($C_i = 1$), because
what is destroyed is the marginal challenge, not the whole stock. Third, the
penalty is \emph{linearly decaying} in the inning: a failed first-inning
challenge forfeits eight innings of option value and is charged
$-(8/9) \times EV_c/C_i$, whereas a failed ninth-inning challenge forfeits
nothing and is charged zero. Fourth, extra innings floor at zero via the
$\max(0,\cdot)$ term, consistent with inning-level challenge resets.

Both the linear decay and the $1/C$ share are modeling choices, not derived
results: they are the simplest forms with the correct boundary behavior, and
they avoid free parameters the current data cannot identify. The full
solution (which the tennis and cricket challenge literature obtains by
dynamic programming \citep{abramitzky2012,clarkenorman2012,shivakumar2018}) would
price the marginal value of the $k$-th challenge and the option value of
challenges held. A convex decay, reflecting that late innings are more
leveraged per inning than early ones, is a plausible refinement.

\subsubsection{Unified metric}
\begin{align}
cRV_i = \Delta RV_{pitch,i} + Penalty_i.
\end{align}
Because exactly one of the two terms is nonzero for any given event
(successes have no penalty, failures have no delta), $cRV$ decomposes cleanly:
positive values come entirely from conversions, negative values entirely from
wasted challenges.

\subsubsection{Terminal outcomes}
\label{sec:terminal}
The central modeling question is how to value challenges that \emph{end the
plate appearance}. A called strike on an 0--0 count merely moves the count; a
called strike on a 3--2 count ends the at-bat. Since challenges are
disproportionately spent on full counts and two-strike counts (exactly where
the stakes are highest), getting terminal outcomes right dominates the metric's
behavior.

We handle three cases:

\paragraph{Walk (Ball~4).} Runners advance under standard force rules: the
batter takes first; a runner on first is forced to second; a runner on second
advances only if first was occupied; a runner on third scores only if the bases
were loaded. Any forced run is added immediately to the state's value. Remaining
run expectancy is then looked up at the post-walk base state, unchanged outs,
and a reset $0$--$0$ count (a new batter is up).

\paragraph{Strikeout (Strike~3).} Outs increment by one. If outs reach three,
the half-inning ends and the remaining run expectancy is zero. Otherwise RE is
looked up at the same bases, outs$+1$, and a reset $0$--$0$ count. A pure
strikeout scores no runs, so the immediate-runs term is zero.

\paragraph{Non-terminal count change.} Bases and outs are unchanged; only the
count state updates in the RE288 lookup. These are the low-magnitude events: an
0--1 to 1--0 swing is worth a few hundredths of a run.

Formal transition rules are given in Appendix~\ref{app:transitions}.

\subsubsection{Why the terminal treatment matters}
Earlier drafts of this pipeline assigned $RE = 0$ to \emph{all} terminal
states, on the reasoning that a completed plate appearance has no
count-conditioned run expectancy. That is a category error: the plate
appearance ends, but the \emph{inning} does not. A bases-loaded walk leaves
the bases loaded with a run already in; a strikeout with one out leaves runners
on with two outs. Both states carry substantial run expectancy.

The practical consequence is severe. Under the old treatment, a bases-loaded
walk conversion scores $\Delta RV = 0$ instead of the $+1.91$ runs it is
actually worth, and an inning-ending strikeout conversion scores $0$ instead of
$+1.72$. On the 2026 terminal-challenge corpus (Section~\ref{sec:results}), zeroing terminal
states drives mean $cRV$ to $-0.016$: the metric reports that challenging is,
on net, \emph{harmful}, because only the penalties survive. With terminal states
valued correctly, mean $cRV$ is $+0.159$. The sign of the headline result flips
on this one modeling decision, and it flips on real data, not just a
constructed example.

\subsection{Data and pipeline}
\label{sec:pipeline}
The repository implements five core modules plus a validation layer
(Figure~\ref{fig:pipeline}):

\begin{enumerate}
  \item \textbf{Ingestion} (\texttt{ingestion.py}): Statcast pitch-by-pitch
        pull via \texttt{pybaseball}, defaulting to the 2026 ABS-era window.
        Returns an empty frame on network failure rather than raising, so the
        pipeline degrades gracefully.
  \item \textbf{Event isolation} (\texttt{events.py}): filters the
        \texttt{des} field for challenge language, infers challenger role
        (batter/catcher), called pitch type, and success/failure, and flags
        parse-quality issues.
  \item \textbf{Alternate realities} (\texttt{realities.py}): from the
        pre-pitch count, bases, and outs, constructs the umpire-call state
        and the overturned state, including walk force chains and
        strikeout out increments.
  \item \textbf{RE288 mapping} (\texttt{re288.py}): loads an empirical matrix
        estimated from play-by-play when available, else a structured
        synthetic surface; maps both realities to run values.
  \item \textbf{Scoring} (\texttt{calculate.py}): applies $\Delta RV$ and the
        opportunity-cost penalty.
\end{enumerate}

\crvfigure{fig:pipeline}{pipeline_diagram}{1.0}{Five-module $cRV$ pipeline: ingestion, event isolation, alternate-state generation, RE288 mapping, and scoring. Each stage writes an inspectable intermediate file, so any challenge event can be traced from its raw description text through to its final score.}

\subsubsection{Event isolation and parse quality}
Challenge events are identified from free-text descriptions requiring both
challenge language and an outcome token (``overturned'' / ``upheld'' /
``stands''). Role and called-pitch-type are inferred from the same text. When
inference fails (for example, a description reading only ``Challenge overturned
to a ball'' with no role named), the event is retained but flagged via a
\texttt{parse\_quality} column rather than silently dropped or silently
mis-assigned.

This matters for honest accounting: in the 2026 corpus, only 2 of $1{,}971$
events carry a \texttt{missing\_challenger} flag. Such events contribute zero to
$\Delta RV$ (no role means no sign convention) and appear as an
\texttt{unknown} row in the role summary. Production ABS feeds with structured
challenge flags would eliminate this ambiguity entirely.

A structural feature of the extraction deserves note. Statcast records the
descriptive \texttt{des} text on the final pitch of a plate appearance, so the
challenges recoverable from it are precisely those that end the plate
appearance: $1{,}413$ overturned or upheld strikeout calls and $558$ walk
calls, concentrated on two-strike and three-ball counts. Non-terminal
count-change challenges (for example an 0--1 call overturned to 1--0 mid
at-bat) are not separately identifiable from \texttt{des} and are absent from
this corpus. This is a selection property of the data source, not of the metric;
Section~\ref{sec:results} reports on the terminal challenges that dominate
actual challenge usage.

\subsubsection{RE288 estimation}
The empirical estimator groups play-by-play pitches by
(base state, outs, count), computes runs scored from that pitch to the end of
the half-inning, and averages. It also reports a per-cell
\texttt{sample\_size}, since rare states (3--0 counts with the bases loaded
and two outs) are noisy and a multi-season baseline is required for stability.
Across the 288 cells the median cell rests on 585 pitches and the mean on
$1{,}858$, but the distribution is right-skewed: 28 cells (10\%) have fewer
than 100 pitches and 188 cells (65\%) fewer than $1{,}000$
(Table~\ref{tab:re288samples}). The sparsest cells are precisely the rare
base-out-count states that most need a multi-season baseline.

\begin{table}[h]
\centering
\caption{Per-cell sample sizes of the empirical RE288 surface.}
\label{tab:re288samples}
\small
\begin{tabular}{lr}
\toprule
Metric & Value \\
\midrule
RE288 cells & $288$ \\
Min pitches per cell & $11$ \\
Median pitches per cell & $585$ \\
Mean pitches per cell & $1858.160$ \\
Max pitches per cell & $33388$ \\
Cells under 100 pitches & $28$ \\
Cells under 1000 pitches & $188$ \\
\bottomrule
\end{tabular}
\end{table}


When no play-by-play is available, a structured synthetic surface is used:
$RE = 0.56 + 0.18r - 0.32o + 0.04b - 0.03s$, where $r$ is runners on base, $o$
is outs, and $b, s$ are balls and strikes. This is monotone in the right
directions and adequate for exercising the pipeline, but it is not a baseball
estimate and results derived from it are labeled accordingly.

\subsubsection{Provenance and reproducibility}
All artifacts regenerate with:
\begin{quote}
\texttt{python pipeline.py --output-dir outputs}\\
\texttt{python pipeline.py --force-synthetic  \# offline demo corpus}
\end{quote}
A \texttt{run\crvus{}metadata.csv} records the data source
(\texttt{statcast}\slash\texttt{synthetic}) and RE288 source
(\texttt{empirical}\slash\texttt{synthetic}\slash\texttt{csv}) for every run, so
no figure can be silently misattributed. The results in this paper were
produced with \texttt{data\crvus{}source = statcast} and
\texttt{re288\crvus{}source = empirical}: real 2026 challenge events scored
against an RE288 matrix estimated from the same season's play-by-play.

\subsection{Validation design}
\label{sec:validation}
The state-transition logic is the part of the pipeline most prone to silent
error, so its correctness is established independently of any particular data
pull, against a \emph{designed} corpus that exercises every transition path:

\begin{itemize}
  \item non-terminal count changes (0--0 strike overturned to a ball),
  \item a bases-loaded walk conversion (run forced in, bases stay loaded),
  \item an inning-ending strikeout conversion (RE drops to zero),
  \item a two-out strikeout that does \emph{not} end the inning,
  \item an extra-inning failure (penalty must be exactly zero),
  \item early and late failures (penalty must scale with innings remaining),
  \item a parse-ambiguous description (role inference must fail loudly).
\end{itemize}

A unit-test suite (\texttt{tests/}) asserts the expected behavior of each path
independently of the pipeline: walk force chains for all eight base states,
terminal RE zeroing, extra-inning penalty nullification, and parser role
extraction. All tests pass on the current implementation. This harness makes the
empirical results of Section~\ref{sec:results} regression-safe: the transition
rules that value a bases-loaded walk are the same rules the tests pin down.

\section{Results}
\label{sec:results}
We apply the pipeline to every plate-appearance-ending ABS challenge
recoverable from the 2026 regular season: $1{,}971$ events spanning
$2026$-$03$-$26$ through $2026$-$08$-$12$,
issued by $446$ distinct challengers, scored against an RE288 matrix estimated
from the same season's play-by-play ($EV_c = 0.15$). All artifacts live in
\texttt{outputs/}; the provenance metadata confirms the \texttt{statcast} and
\texttt{empirical} sources.

\begin{table}[h]
\centering
\caption{Validation summary metrics for the current run.}
\label{tab:validation}
\small
\begin{tabular}{lrrrl}
\toprule
Metric & Value & CI lower & CI upper & Notes \\
\midrule
Total events & $1914.000$ & $1914.000$ & $1914.000$ & Count of challenge events \\
Success rate & $0.460$ & $0.000$ & $1.000$ & Share of successful challenges \\
Mean cRV & $0.143$ & $0.130$ & $0.158$ & Bootstrap CI for mean cRV \\
\bottomrule
\end{tabular}
\end{table}


Mean $cRV$ is $0.159$ runs per challenge, with a game-clustered bootstrap
$95\%$ CI of $[0.147,\,0.172]$, comfortably above zero, so the challenge as
actually used in 2026 added run value on net within this terminal-challenge
census. A positive mean is the expected-value answer for the challenge system,
not a scoring artifact: a successful overturn realizes a positive run swing,
while a failure is charged only the opportunity cost of the challenge it
destroyed (Section~\ref{sec:limitations}). The overall success rate is
$46.5\%$. At $n = 1{,}971$ the interval is tight, a marked contrast to the wide
bands a small validation corpus would produce, and it is narrow enough to
support role- and leverage-level comparisons rather than only a directional
claim.

\begin{table}[h]
\centering
\caption{Top player-level cumulative cRV values.}
\label{tab:leaderboard}
\small
\begin{tabular}{llrrrr}
\toprule
Player & Role & Challenges & Success rate & Cumulative $cRV$ & Mean $cRV$ \\
\midrule
Shea Langeliers & catcher & 33 & $0.697$ & $8.791$ & $0.266$ \\
Edgar Quero & catcher & 18 & $0.667$ & $7.031$ & $0.391$ \\
Hunter Goodman & catcher & 26 & $0.692$ & $6.481$ & $0.249$ \\
Carson Kelly & catcher & 17 & $0.882$ & $5.483$ & $0.323$ \\
Victor Caratini & catcher & 24 & $0.792$ & $5.433$ & $0.226$ \\
Freddy Fermin & catcher & 27 & $0.667$ & $5.422$ & $0.201$ \\
Ryan Jeffers & catcher & 25 & $0.640$ & $5.168$ & $0.207$ \\
Jonah Heim & catcher & 21 & $0.762$ & $5.118$ & $0.244$ \\
Francisco Alvarez & catcher & 18 & $0.667$ & $5.095$ & $0.283$ \\
Carlos Narváez & catcher & 18 & $0.889$ & $5.069$ & $0.282$ \\
\bottomrule
\end{tabular}
\end{table}


\begin{table}[h]
\centering
\caption{Aggregate cRV by challenger role.}
\label{tab:challenger}
\small
\begin{tabular}{lrrr}
\toprule
Role & Challenges & Cumulative $cRV$ & Mean $cRV$ \\
\midrule
catcher & 952 & $211.657$ & $0.222$ \\
batter & 960 & $62.782$ & $0.065$ \\
unknown & 2 & $-0.033$ & $-0.017$ \\
\bottomrule
\end{tabular}
\end{table}


Figure~\ref{fig:leaderboard} shows the player-level distribution and
Figure~\ref{fig:challenger} the same totals aggregated by role. The shape is the
one the framework predicts: value is concentrated in terminal-state conversions,
while failures contribute a thin band of small negative values whose magnitude
is set entirely by inning. The role split is the headline empirical finding.
Batters and catchers issue challenges in almost equal numbers ($985$ versus
$984$), but catchers accumulate $227.6$ runs of $cRV$ to the batters' $86.7$, a
per-challenge mean of $0.231$ against $0.088$. This comparison is
descriptive, not causal: catchers and batters issue challenges in structurally
different spots (converting two-strike takes into outs versus avoiding called
third strikes), so it does not by itself identify superior judgment. Catchers
convert borderline two-strike takes into outs, the single most valuable thing a
challenge can do; batters more often convert a called third strike into a mere
continuation of the at-bat.

\crvfigure{fig:leaderboard}{leaderboard_top10}{0.9}{Player-level cumulative $cRV$ for the top challengers of the 2026 regular season. Bars extending right of the zero rule (blue) are positive and arise entirely from successful conversions; bars extending left (red) are negative and arise entirely from opportunity-cost penalties on failed challenges. Bar direction relative to the zero rule encodes the sign independently of color, so the comparison survives grayscale reproduction. The leaderboard is dominated by catchers, the primary users of the challenge.}

\crvfigure{fig:challenger}{challenger_totals}{0.8}{Cumulative $cRV$ aggregated by challenger role (batter, catcher, and a small number of events whose role could not be parsed). Batters and catchers issue challenges in nearly equal numbers, but catchers account for the large majority of aggregate value, converting borderline two-strike calls into inning-ending outs.}

\subsection{Information and reliability}
Two checks bear on whether $cRV$ measures anything beyond raw accuracy, and
whether it is stable enough to be a skill rather than noise
(Table~\ref{tab:info}). First, across the $47$ players with at least $10$
challenges, mean $cRV$ correlates $0.72$ (Pearson) and $0.68$ (Spearman) with
challenge success rate. Accuracy explains roughly half the variance, but the
remaining half is leverage and timing, the information the run-denominated
metric adds over a success-rate ranking. Second, split-half reliability within
2026 is weaker: for the $45$ players with at least four challenges in each half
of the season, mean $cRV$ in the first half correlates $0.34$ (Pearson) with
the second half. Player-level $cRV$ is therefore only modestly stable within a
half-season, consistent with a young, sparse sample; the run-denominated
ordering is robust in aggregate but not yet a per-player skill signal.

\begin{table}[h]
\centering
\caption{Information and reliability of player-level cRV.}
\label{tab:info}
\small
\begin{tabular}{lr}
\toprule
Metric & Value \\
\midrule
Qualifying players & $47$ \\
Min challenges per player & $10$ \\
Pearson r (cRV vs success rate) & $0.723$ \\
Spearman rho (cRV vs success rate) & $0.678$ \\
Qualifying players & $45$ \\
Min challenges per half & $4$ \\
Pearson r (split-half cRV) & $0.344$ \\
Spearman rho (split-half cRV) & $0.333$ \\
\bottomrule
\end{tabular}
\end{table}


\subsection{Case studies}
Three events illustrate the model's economic content. Each is traced from
pre-pitch state to final score.

\paragraph{Bases-loaded walk conversion (Eli White, batter, $+1.91$).}
Fifth inning, bases loaded, 1 out, 3--2 count. The umpire calls strike three;
White challenges and wins, converting the pitch to ball four.
\begin{itemize}
  \item \emph{Original reality:} strikeout. Outs $1 \to 2$, bases loaded
        unchanged. $RE(\text{original}) = 0.721$.
  \item \emph{Overturned reality:} walk. The runner on third is forced home
        ($+1$ run, Ozzie Albies scores), bases remain loaded, outs stay at 1.
        $RE(\text{overturned}) = 2.633$.
\end{itemize}
Batter delta: $2.633 - 0.721 = 1.912$ runs, with no penalty (the challenge is
retained). A forced run plus the gap between one and two outs with the bases
loaded is an enormous swing, and it is the single most valuable challenge
outcome available.

\paragraph{Inning-ending strikeout conversion (Luis Campusano, catcher, $+1.72$).}
Fourth inning, bases loaded, 2 outs, 3--2 count. The umpire calls ball four;
Campusano challenges and wins, converting the walk into strike three.
\begin{itemize}
  \item \emph{Original reality:} walk. A run is forced in ($+1$), bases stay
        loaded, still 2 outs. $RE(\text{original}) = 1.721$.
  \item \emph{Overturned reality:} strikeout, third out. Inning over.
        $RE(\text{overturned}) = 0$.
\end{itemize}
Catcher delta: $1.721 - 0 = 1.721$ runs. This is exactly the case the old
zero-everything treatment got most wrong: it scored the inning-ending
conversion (erasing a run already in and ending the threat) at zero.

\paragraph{Early failed challenge (Carter Jensen, catcher, $-0.07$).}
A first-inning failed catcher challenge (the team's first failure, so two
challenges were still held) yields $\Delta RV = 0$ and
$Penalty = -(8/9) \times 0.15 \times (1/2) \approx -0.067$. Had it been the
team's last challenge the charge would double to $-0.133$, and any ninth-inning
failure yields exactly $0.00$ because no innings of option value remain. Same
outcome and accuracy, but the price depends on both how much game remains and
how much of the challenge budget survives, the resource-management incentive
the framework exists to encode.

\subsection{Leverage stratification}
\begin{table}[h]
\centering
\caption{cRV by inning leverage bucket.}
\label{tab:leverage}
\small
\begin{tabular}{lrrr}
\toprule
Leverage bucket & Events & Mean $cRV$ & Total $cRV$ \\
\midrule
early(1-3) & 557 & $0.138$ & $76.802$ \\
mid(4-6) & 617 & $0.153$ & $94.291$ \\
late(7-9) & 713 & $0.140$ & $99.776$ \\
extra(10+) & 27 & $0.137$ & $3.686$ \\
\bottomrule
\end{tabular}
\end{table}


Total value tracks event count: the late bucket (innings 7--9) carries the most
challenges ($735$, $105.6$ runs), followed by mid ($631$ events, $108.0$) and
early ($578$ events, $97.0$), with a small extra-inning tail ($27$ events,
$3.8$). The \emph{mean} is mildly higher early and mid-game than late: $0.168$,
$0.171$, $0.144$, and $0.139$ runs respectively. Two forces pull against each
other: the inning decay makes an early failure costlier, but the scarcity
weight makes a team's \emph{first} failure (disproportionately early) only
half as costly as its last. Terminal conversions dominate the numerator in
every inning, so per-challenge value stays within a narrow band while total
value accumulates wherever challenges are most frequent.

\subsection{Sensitivity to $EV_c$}
\label{sec:sensitivity}
\begin{table}[h]
\centering
\caption{Sensitivity of aggregate cRV to the baseline challenge value $EV_c$.}
\label{tab:evsens}
\small
\begin{tabular}{rrr}
\toprule
$EV_c$ & Mean $cRV$ & Total $cRV$ \\
\midrule
$0.050$ & $0.164$ & $314.388$ \\
$0.100$ & $0.154$ & $294.472$ \\
$0.150$ & $0.143$ & $274.555$ \\
$0.200$ & $0.133$ & $254.638$ \\
$0.250$ & $0.123$ & $234.722$ \\
$0.300$ & $0.112$ & $214.805$ \\
\bottomrule
\end{tabular}
\end{table}


\crvfigure{fig:sensitivity}{ev_sensitivity}{0.85}{Total corpus $cRV$ as a function of the baseline retained-challenge value $EV_c$, over the range $[0.034, 0.30]$ runs spanning the empirical estimate and the fixed baseline. The relationship is exactly linear because the opportunity-cost penalty applies only to failed challenges and scales linearly in $EV_c$; successful challenges carry no penalty term. Endpoint values are labeled.}

Total corpus $cRV$ ranges from $338.9$ runs at the empirical $EV_c = 0.034$ to
$282.6$ runs at $EV_c = 0.30$ (Figure~\ref{fig:sensitivity}), a spread of
$56.2$ runs across the grid, with the mean falling from $0.172$ to $0.143$.
The relationship is exactly linear: from Equation~\ref{eq:penalty},
$\sum_i cRV_i = \sum_i \Delta RV_i - EV_c \sum_{i \in \text{fail}}
(9 - \mathrm{inning}_i)/(9\,C_i)$, and only the second term depends on $EV_c$.

Two conclusions follow. First, because $cRV$ is monotone nonincreasing in
$EV_c$ (only the penalty depends on it), the positive sign of mean $cRV$ holds
across the entire grid (at the data-driven value $0.034$ as well as the fixed
baseline), so the conclusion that challenging adds value is not an artifact of
the $EV_c$ choice. Second, the \emph{ranking} of successful
high-leverage events is completely invariant to $EV_c$, since successes carry
no penalty term at all. $EV_c$ affects only how harshly failures are judged
relative to one another, and even there it is a monotone rescaling.

\subsection{Win-probability overlay}
As a win-denominated complement to the run-denominated results, we define a
descriptive win-denominated analog of $cRV$: a successful overturn is priced by its
win-probability swing, and a failed challenge carries a win-denominated
opportunity-cost penalty. Statcast records
\texttt{delta\crvus{}home\crvus{}win\crvus{}exp}, the change in the home team's
win probability on the resolved pitch; on a successful challenge that pitch is
the overturned call, so signing it to the challenging side measures the win
probability the overturn actually moved. The penalty mirrors
Equation~\ref{eq:penalty}, with the retained-challenge value now measured in
win probability: $0.0025$ win probability per retained challenge, estimated
from the same three-factor decomposition as the run version
(Table~\ref{tab:evcwin}). The resulting cWPA preserves the run-denominated
ordering: catchers accumulate $+14.9$ win probability ($0.0151$ per challenge)
to batters' $+5.3$ ($0.0054$ per challenge; Table~\ref{tab:cwpa}), the same
catcher dominance found in runs, for a corpus total of $+20.2$ win probability
(Figure~\ref{fig:cwpa}).

\begin{table}[h]
\centering
\caption{Empirical estimate of the win-denominated retained-challenge value.}
\label{tab:evcwin}
\small
\begin{tabular}{lrrr}
\toprule
Metric & Value & CI lower & CI upper \\
\midrule
P(future need) & $0.206$ & $0.191$ & $0.222$ \\
P(win) & $0.465$ & $0.444$ & $0.486$ \\
Mean late-game win swing & $0.026$ & $0.023$ & $0.029$ \\
Win-denominated EV\_c & $0.002$ & $0.002$ & $0.003$ \\
\bottomrule
\end{tabular}
\end{table}


\begin{table}[h]
\centering
\caption{Realized win-probability swing of successful overturns by role.}
\label{tab:cwpa}
\small
\begin{tabular}{lrrr}
\toprule
Role & Events & Total cWPA & Mean cWPA \\
\midrule
catcher & 984 & $14.895$ & $0.015$ \\
batter & 985 & $5.324$ & $0.005$ \\
unknown & 2 & $0.000$ & $0.000$ \\
\bottomrule
\end{tabular}
\end{table}


\crvfigure{fig:cwpa}{challenge_wpa}{0.8}{Total challenge win probability (cWPA) by challenger role, the win-denominated counterpart of Figure~3. Successful overturns are priced by their win-probability swing and failed challenges by a win-denominated opportunity-cost penalty. As in runs, catchers accumulate the large majority of win-probability value.}

\section{Discussion}
\subsection{Managerial implications}
Under this framework, challenging a borderline strike on an 0--0 count is
rarely justified. The count move (0--1 $\to$ 1--0) is worth roughly $0.07$
runs in our surface, while a first-inning miss costs $-(8/9) \times 0.15 / C$,
about $-0.07$ runs if it is the team's first failure ($C = 2$) and $-0.13$ if
it is the last ($C = 1$) under the baseline $EV_c$. Break-even therefore
requires a win probability of roughly $49\%$ to $65\%$, depending on how much
of the challenge budget remains, which is higher than typical challenge success rates.
Under the empirical $EV_c = 0.034$ those first-inning penalties shrink to
$-0.015$ and $-0.03$ runs, and break-even falls to about $18\%$--$30\%$, so
the strength of this recommendation turns on which $EV_c$ one trusts.

By contrast, full-count terminal challenges (especially with runners on and
two outs) can swing between half a run and a run and a half of expectancy.
At those magnitudes, break-even win probability falls into the low single
digits of percent. Players should be instructed to spend challenges almost
freely on plate-appearance-ending calls with runners in scoring position, and
almost never on early-count calls with the bases empty.

More generally, cumulative $cRV$ is a decision-quality signal distinct from raw
challenge accuracy. A catcher who goes 5-for-10 on low-leverage challenges may
post a worse $cRV$ than one who goes 3-for-5 exclusively in high-leverage
spots. Accuracy measures the eye; $cRV$ measures the judgment.

\subsection{Game-theoretic spillovers}
Once a team has exhausted its challenges, the strategic environment changes for
both sides. Pitchers may expand the zone and catchers may set up further off
the plate, knowing borderline takes cannot be reviewed. Conversely, a team
holding challenges late exerts a deterrent effect that never appears in any
box score.

$cRV$ as currently specified prices neither effect. It values challenges
actually thrown, not the option value of challenges held. A complete
game-theoretic treatment would model the equilibrium zone expansion as a
function of remaining challenges on both sides, a substantially harder problem
requiring pitch-location data and a behavioral model of umpire and pitcher
response.

\subsection{Relationship to framing}
In a full-ABS-challenge world, framing runs should compress toward zero for
challenged pitches while remaining live for unchallenged ones. The two metrics
therefore partition the borderline-pitch value space: framing governs the
pitches nobody reviews, $cRV$ governs the ones somebody does. A complete
catcher valuation in the ABS era plausibly sums the two, though care is needed
to avoid double-counting pitches that were framed \emph{and} challenged.

\subsection{Limitations}
\label{sec:limitations}
\begin{itemize}
  \item \textbf{Data source and coverage.} The corpus is the full set of
        2026 regular-season challenge events recoverable from Statcast
        play-by-play descriptions ($n = 1{,}971$). Because Statcast records an
        outcome in the free-text description only on the pitch that ends the
        plate appearance, every recovered challenge is plate-appearance-ending
        ($1{,}413$ strikeout calls, $558$ walk calls). Count-changing
        challenges that leave the plate appearance live are not represented.
        The results are therefore a census of \emph{terminal} challenges rather
        than a sample of all challenges, which captures the highest-leverage
        events but omits the low-value count-changing tail, so the reported
        mean is upward-biased relative to the mean over all challenges.
  \item \textbf{RE288 surface.} The matrix is estimated empirically from 2026
        regular-season play-by-play
        (\texttt{estimate\crvus{}empirical\crvus{}re288\crvus{}matrix}); cells
        visited rarely inherit correspondingly wide sampling error, and
        multi-season pooling would stabilize them.
  \item \textbf{Sample size and selection.} $n = 1{,}971$ is a season-to-date
        census of terminal challenges, not a random sample. The bootstrap
        interval $[0.147,\,0.172]$ reflects resampling variability within that
        census, not sampling from a larger population, and the reported mean is
        an average over the highest-leverage subset of challenges.
  \item \textbf{Opportunity-cost specification.} The retained-challenge
        penalty (Equation~\ref{eq:penalty}) conditions on innings remaining and
        on the number of challenges held before a failure (a $1/C$ share), but
        it is a first-order approximation, not the full dynamic program solved
        in the tennis and cricket challenge literature
        \citep{abramitzky2012,clarkenorman2012,shivakumar2018}. A positive mean
        $cRV$ is nonetheless the expected-value answer (a successful overturn
        realizes a positive run swing while a failure is charged the
        opportunity cost of the challenge it destroyed), rather than an
        artifact of an asymmetric scoring rule.
  \item \textbf{Fixed $EV_c$.} The baseline challenge value is fixed at
        $0.15$ runs for the headline results. The empirical estimate ($0.034$
        runs) is about a quarter of that, so the headline results are
        conservative; Section~\ref{sec:sensitivity} bounds the model risk
        across both values and beyond.
  \item \textbf{Decay and share assumptions.} The penalty decays linearly in
        innings remaining and divides by the number of challenges held. Both
        have correct boundary behavior but no empirical justification; a convex
        decay, or a data-driven marginal value for the $k$-th challenge, may
        better reflect late-inning leverage and option value.
  \item \textbf{Parse ambiguity.} Role inference from free-text descriptions
        can fail ($2$ of $1{,}971$ events). Production feeds with
        structured challenge flags would eliminate this.
  \item \textbf{Win-probability layer is descriptive.} $cRV$ is
        run-denominated, not win-denominated, so a one-run swing in a blowout
        and in a tie game score identically. The win-denominated overlay prices
        overturns and their penalty in win probability, but it leans on
        Statcast's pitch-level win-expectancy swings rather than a full
        leverage-index model of challenge value.
  \item \textbf{Challenges held are partially priced.} The penalty charges the
        $1/C$ share of the challenge a failure destroys, but it does not price
        the option value of challenges retained for future use, nor the
        deterrent effect a holding team exerts on an opponent.
\end{itemize}

\subsection{Future work}
\label{sec:future}
Immediate priorities, roughly in dependency order:
\begin{enumerate}
  \item Recover count-changing (non-terminal) challenges that Statcast's
        free-text descriptions omit, using structured ABS challenge logs, so the
        corpus becomes a census of \emph{all} challenges rather than terminal
        ones only.
  \item Build a multi-season empirical RE288 matrix and report per-cell sample
        sizes, flagging low-$n$ states as unstable.
  \item Validate description parsers against official ABS game logs and
        quantify false-positive and false-negative rates.
  \item Replace the static three-factor $EV_c$ estimate and the $1/C$ share
        with the full dynamic program of the tennis and cricket challenge
        literature \citep{abramitzky2012,clarkenorman2012,shivakumar2018},
        which prices the marginal value of the $k$-th challenge and the option
        value of challenges held.
  \item Extend the within-season split-half reliability check (currently
        $r \approx 0.34$) across multiple seasons and larger per-player
        samples, to establish whether player $cRV$ stabilizes into a skill
        signal.
  \item Explore dynamic $P(\text{win})$ conditioned on pitch location and
        historical umpire/ABS agreement rates, enabling a prospective
        ``should you challenge this?'' model rather than a retrospective score.
  \item Layer a full leverage-index model over $cRV$, conditioning the win
        value of an overturn on the game state rather than Statcast's average
        pitch-level win-expectancy swing.
\end{enumerate}

\section{Conclusion}
$cRV$ provides a practical, run-based framework for measuring challenge
strategy in the ABS era. By pairing RE288 state deltas (including correct walk
and strikeout transitions) with a transparent opportunity-cost penalty, the
metric distinguishes high-leverage conversion skill from reckless early usage,
a distinction invisible to raw success rates.

The methodological finding we would most emphasize is the terminal-outcome
treatment. Plate-appearance-ending challenges are simultaneously the most
common high-leverage challenges and the easiest to model incorrectly; on the
2026 terminal-challenge corpus, zeroing them turns a mean $cRV$ of $+0.159$ into $-0.016$,
inverting the sign of the headline result. Any future implementation of a
challenge-valuation metric should treat walk force chains and strikeout out
increments as first-class cases, not edge cases.

The accompanying open pipeline makes every table and figure in this paper
regenerable from a single command, with provenance tracking that records
whether each run drew on live Statcast data or the synthetic validation corpus.
Applied to the $1{,}971$ plate-appearance-ending challenges recoverable from
the 2026 season, it establishes catchers as the
primary beneficiaries of the ABS challenge and terminal conversions as the
events that carry the metric.

\appendix
\section{State-transition rules}
\label{app:transitions}

Let the pre-pitch state be count $(b, s)$, base string
$\beta = \beta_1\beta_2\beta_3$, and outs $o$.

\subsection{Called strike}
\begin{align}
(b, s) \to
\begin{cases}
(b, s+1), \; \beta, \; o & \text{if } s + 1 < 3 \quad \text{(non-terminal)},\\
(0,0), \; \beta, \; o+1 & \text{if } s + 1 = 3, \; o + 1 < 3
  \quad \text{(strikeout)},\\
\text{inning over}, \; RE = 0 & \text{if } s + 1 = 3, \; o + 1 = 3.
\end{cases}
\end{align}
Runs scored on the play: $0$ in all cases.

\subsection{Called ball}
\begin{align}
(b, s) \to
\begin{cases}
(b+1, s), \; \beta, \; o & \text{if } b + 1 < 4 \quad \text{(non-terminal)},\\
(0,0), \; \beta', \; o & \text{if } b + 1 = 4 \quad \text{(walk)},
\end{cases}
\end{align}
where the post-walk base string $\beta'$ and runs scored $\rho$ follow the
force chain:
\begin{align}
\beta'_1 &= 1 \quad \text{(batter to first, always)},\\
\beta'_2 &= \begin{cases} 1 & \text{if } \beta_1 = 1 \text{ (forced)},\\
                          \beta_2 & \text{otherwise},\end{cases}\\
\beta'_3 &= \begin{cases} 1 & \text{if } \beta_1 = \beta_2 = 1
                              \text{ (forced from second)},\\
                          \beta_3 & \text{otherwise},\end{cases}\\
\rho &= \begin{cases} 1 & \text{if } \beta = 111 \text{ (bases loaded)},\\
                      0 & \text{otherwise}.\end{cases}
\end{align}

Enumerating all eight base states:

\begin{table}[h]
\centering
\caption{Walk force chain for all base states.}
\label{tab:forcechain}
\begin{tabular}{llll}
\toprule
Pre-walk $\beta$ & Description & Post-walk $\beta'$ & Runs $\rho$ \\
\midrule
000 & empty          & 100 & 0 \\
100 & 1B             & 110 & 0 \\
010 & 2B             & 110 & 0 \\
001 & 3B             & 101 & 0 \\
110 & 1B, 2B         & 111 & 0 \\
101 & 1B, 3B         & 111 & 0 \\
011 & 2B, 3B         & 111 & 0 \\
111 & loaded         & 111 & 1 \\
\bottomrule
\end{tabular}
\end{table}

Note the unforced cases: with a runner on second only ($010$), the batter takes
first and the runner \emph{holds}, giving $110$. With runners on second and
third ($011$), both hold and the batter takes first, giving $111$ with no run
scored; the runner on third is not forced.

\subsection{Overturn mapping}
An overturned called strike is evaluated as a called ball, and an overturned
called ball as a called strike. If the challenge fails, the overturned reality
is set equal to the original reality, guaranteeing $\Delta RV = 0$ by
construction rather than by a separate branch.

\subsection{State value}
\begin{align}
RE(\text{state}) =
\begin{cases}
\rho & \text{if inning over or } o \geq 3,\\
\rho + RE288(\beta, o, \text{count}) & \text{otherwise}.
\end{cases}
\end{align}

\section{Reproducing the results}
\label{app:repro}
\begin{quote}
\begin{verbatim}
python3 -m venv .venv && source .venv/bin/activate
pip install -r requirements.txt
python pipeline.py --output-dir outputs        # full 2026 Statcast pull
python -m pytest tests/ -q
cd paper && make paper
\end{verbatim}
\end{quote}
The first \texttt{pipeline.py} invocation fetches and caches the 2026
regular-season Statcast feed, isolates challenge events, estimates the empirical
RE288 surface, and writes every table and figure; add
\texttt{-{}-force-synthetic} to regenerate the offline validation corpus instead.
Every table in Section~\ref{sec:results} is generated by the pipeline and
included via \verb|\input|; rerunning the pipeline updates the manuscript
automatically. A \texttt{run\_metadata.csv} artifact records the data
provenance for the run that produced the current tables. The analysis code and
the data-generating scripts will be deposited in a public repository on
acceptance; the repository is omitted here to preserve blind review.



\section*{Data Availability}
The pitch-level inputs are derived from publicly available Statcast data,
retrievable through the \texttt{pybaseball} interface \citep{pybaseball}. The
results reported in this manuscript use the designed synthetic validation
corpus described in Section~\ref{sec:validation}, which is generated
deterministically by the accompanying code.

\begin{acknowledgement}
  TODO(post-acceptance): restore acknowledgments here.
\end{acknowledgement}

\begin{funding}
  No external funding supported this work.
\end{funding}

\bibliography{references}

\end{document}
